% Conditions that often produce no symptoms; a clinical-state index, not a disease category.
\small
\noindent\textit{Scope}: This reference list groups diseases, infections, and other conditions that may be present without noticeable symptoms, particularly during early, latent, or compensated stages. Being asymptomatic is a clinical state or mode of presentation rather than a separate category of disease, so every condition here also belongs under its usual organ-system or specialty heading, and ``May be asymptomatic'' is best treated as an attribute or tag in addition to that classification. Many of these conditions cause symptoms in some people and none in others, and some come to attention only through screening or when a complication appears. The list is a curated, high-importance reference rather than a prevalence ranking, and some conditions appear under more than one heading where they are clinically relevant to several systems. Individual symptoms, signs, and laboratory findings are defined as dictionary entries rather than listed here.

\medskip
\noindent\textbf{Important distinction:} ``Asymptomatic disease'' is not a distinct disease; it describes how a disease presents. Hypertension is a cardiovascular disease that is often asymptomatic, chlamydia is an infection that is frequently asymptomatic, osteoporosis is a skeletal disease that is typically asymptomatic until a fracture, and open-angle glaucoma is an eye disease that is often asymptomatic in its early stages.

\medskip
\begin{description}
\item[Cardiovascular] Hypertension; Atherosclerosis; Coronary Artery Disease; Peripheral Artery Disease; Some Arrhythmias; Early Valvular Heart Disease

\item[Metabolic \& Endocrine] Prediabetes; Early Type 2 Diabetes; Dyslipidemia; Subclinical Hypothyroidism; Subclinical Hyperthyroidism; Hyperuricemia

\item[Kidney \& Urinary] Early Chronic Kidney Disease; Diabetic Kidney Disease; Mild Proteinuria; Microscopic Hematuria; Some Kidney Stones

\item[Liver] Metabolic Dysfunction-Associated Steatotic Liver Disease (MASLD); Chronic Hepatitis B; Chronic Hepatitis C; Early Cirrhosis; Hemochromatosis

\item[Infectious] HIV Infection During Clinical Latency; Hepatitis B; Hepatitis C; Chlamydia; Gonorrhea; HPV Infection; Latent Tuberculosis Infection; Some COVID-19 Infections

\item[Cancer \& Precancer] Early Colorectal Cancer; Cervical Precancer; Early Cervical Cancer; Early Prostate Cancer; Early Kidney Cancer; Some Thyroid Cancers; Early Breast Cancer

\item[Gynecological] HPV Infection; Some Ovarian Cysts; Uterine Fibroids; Early Cervical Abnormalities; Some Sexually Transmitted Infections

\item[Male Reproductive] Early Prostate Cancer; Some Sexually Transmitted Infections; Varicocele; Some Cases of Benign Prostatic Enlargement

\item[Blood] Mild Anemia; Monoclonal Gammopathy of Undetermined Significance (MGUS); Early Chronic Leukemia; Mild Thrombocytopenia

\item[Bone] Osteoporosis; Osteopenia

\item[Eye] Open-Angle Glaucoma; Early Diabetic Retinopathy; Early Age-Related Macular Degeneration

\item[Neurological \& Vascular] Silent Cerebral Infarction; Asymptomatic Carotid Stenosis; Some Unruptured Intracranial Aneurysms

\item[Gastrointestinal] Diverticulosis; Colon Polyps; Some Gallstones; Some Hiatal Hernias

\item[Gallbladder] Asymptomatic Gallstones (``Silent Gallstones''); Gallbladder Polyps

\item[Genetic] Hemochromatosis Before Organ Damage; Some Inherited Thrombophilias; Some Genetic Carrier States

\item[Parasitic \& Latent Infections] Some Malaria Infections; Intestinal Parasitic Infections; Toxoplasmosis and Other Latent Infections
\end{description}
